\section{Exact power counts at selectable phases}\label{sec:phase-counts}

Fix an integer base \(b\ge2\) and a real algebraic irrational \(\alpha\).
For a root length \(h\) and integer \(s\ge2\) put
\[
 B=b^h,\qquad D=B-1,\qquad M=B^s,\qquad L=(M-1)/D.
\]
For an integer \(u>0\) let \(u_\perp\) denote its largest divisor coprime
to \(b\). Ridout's theorem gives, for every \(\epsilon>0\),
\begin{equation}\label{eq:phase-ridout}
 \|k b^m\alpha\|\gg_{\alpha,b,\epsilon}
               b^{-\epsilon m}|k|^{-1-\epsilon}
 \qquad (k\in\Z\setminus\{0\},\ m\ge0).
\end{equation}
Indeed the version with restricted denominator factors gives
\(q_\perp\|q\alpha\|\gg_{\alpha,b,\epsilon}q^{-\epsilon}\).
If the nearest numerator and \(q\) have a common factor, reducing the
fraction only decreases the product on the left
\cite{RidoutPadic}. Take \(q=|k|b^m\) and use \(q_\perp\le |k|\).

\subsection{Simultaneous exclusion of small-residue frequencies}

\begin{lemma}\label{lem:phase-residue}
Let \(s\ge2\) and \(C_0\ge0\) be fixed. Suppose
\(\kappa,\gamma_0,\rho>0\), \(\omega=\kappa+\gamma_0\),
\(2(\omega+\rho)<1\), and \(\rho<1/16\). For every sufficiently large
\(h\), every integer window
\([m_0,m_0+(2s+16)h]\), \(0\le m_0\le C_0h\), contains a run of
\(\gg h\) consecutive phases \(m\) at which the following system has
no solution:
\begin{equation}\label{eq:phase-bad}
 0<|k|\le B^{s+\kappa},\quad |r|\le B^\rho,\quad k\equiv r\pmod D,
 \qquad \|k b^m\alpha\|\le B^{-s+\gamma_0}.
\end{equation}
The excluded residues include \(r=0\).
\end{lemma}

\begin{proof}
Choose \(\sigma>0\) small and \(\Lambda\) so that
\[
 \sigma<1/16,\quad 2\omega+3\sigma<1,\quad
 \omega+\rho+3\sigma<\Lambda<1-\omega-\rho.
\]
Set \(G=\lfloor\sigma h/10\rfloor\), \(H=\lfloor\Lambda h\rfloor\),
\(d_1=2\omega+2\sigma\) and \(d_2=\omega+\rho+2\sigma\).
If there is no run of \(G\) good phases, choose one witness to
\eqref{eq:phase-bad} at each bad phase. Every interval of \(G\)
consecutive phases contains a witness. At relative offset \(j\) put
\[
 \alpha_0=b^{m_0}\alpha,\qquad Q_j=b^j k_j,\qquad
 R_j=b^j r_j,\qquad e_j=Q_j\alpha_0-P_j
\]
with \(k_j>0\) after a change of sign and \(P_j\) the nearest integer.
By \eqref{eq:phase-ridout}, with \(\epsilon\) small in terms of
\(C_0,s,\kappa,\sigma\), all the witnesses satisfy
\begin{equation}\label{eq:phase-range}
 k_j>B^{s-\gamma_0-\sigma},\qquad
 |e_j|\gg B^{-s-\kappa-\sigma/4}.
\end{equation}
The same bounds apply if the relative origin is reset at a subwindow.
In fact the only arithmetic inputs in the proof are
\eqref{eq:phase-range} and the lower bound on \(U+V\alpha_0\)
used below.

For two witnesses \(i<j\) with \(j-i\ge d_1h\), their
\((Q,P)\) projections cannot be proportional. Otherwise
\[
 |e_i|=(Q_i/Q_j)|e_j|
 \le B^{-(j-i)/h+\omega+\sigma}B^{-s+\gamma_0}
 \le B^{-s-\kappa-\sigma}.
\]
For nonzero residues, the full rows \((Q,P,R)\) cannot be proportional
if \(j-i\ge d_2h\). Proportionality would also imply
\(k_i/k_j=r_i/r_j\le B^\rho\), giving the stronger estimate
\(|e_i|\le B^{-s-\kappa-2\sigma}\).

For three witnesses spanning at most \(H\) positions let \(j_*\) be the
latest offset. The integer determinant of the rows
\((k_j,b^{j_*-j}P_j,r_j)\) is divisible by \(D\). Subtracting
\(b^{j_*}\alpha_0\) times the first column from the second bounds its
absolute value by \(6B^{\omega+\rho}b^H<D\) for large \(h\).
The determinant is therefore zero. Scaling rows and columns identifies
its rank with that of the original rows.

First consider a subwindow of length \((s+6)h+1\) with witnesses
throughout at the stated \(G\)-spacing, all of nonzero residue. Reset
its relative origin. Every \(H+1\) consecutive positions and every
overlap between two consecutive such intervals contains two witnesses
at distance at least \(d_2h\), by \(H-2G>d_2h\) with a fixed positive
margin. The rows consequently lie in a single rational plane
\(UQ+VP+WR=0\). Use a cross product from the first such window so that
\(|V|\le 2B^{s+\kappa+\rho+2\Lambda}\). The first and last witnesses
have independent \((Q,P)\) projections, so \(W\ne0\). We have
\begin{equation}\label{eq:plane-phase}
 (U+V\alpha_0)k_j+Wr_j=V b^{-j} e_j .
\end{equation}
If \(V\ne0\), the Ridout bound yields
\(\left|U+V\alpha_0\right|\gg B^{-\sigma/4}/|V|\)
by reducing \(\epsilon\) further. Using
\(2\kappa+2\rho+4\Lambda+\gamma_0<4\) gives
\[
 \frac{|V|B^{-s+\gamma_0}}{|U+V\alpha_0|}
    \ll B^{s+4+\sigma/4}.
\]
For \(j\ge(s+5)h\), all the ratios \(k_j/r_j\) are within
\(O(B^{-1+\sigma/4})=o(B^{-2\rho})\) of
\(-W/(U+V\alpha_0)\). Distinct ratios with denominators at most
\(B^\rho\) differ by at least \(B^{-2\rho}\).
Hence the ratios are equal, and the plane equation makes the
\((Q,P)\) projections proportional in this last segment. The segment
contains two witnesses at distance at least \(d_1h\), a contradiction.

If \(V=0\), the nonzero ratio \(r_j/k_j\) is fixed. Its reduced
denominator \(q\) divides each \(k_j\) and is greater than
\(B^{s-\gamma_0-\sigma-\rho}\). For adjacent witnesses of gap
\(g\le G\) the integer
\[
 k_iP_j-b^g k_jP_i=b^g k_j e_i-k_i e_j
\]
is divisible by \(q\) and has magnitude at most
\((b^G+1)B^\omega<q\). Hence it vanishes, giving proportional
projections throughout the subwindow, again a contradiction.

Return to the full window. For consecutive zero-residue witnesses
\(i<j\), if \(j-i\le H\) the same two-row determinant is divisible by
\(D\) and is smaller than \(D\). Their projections are proportional.
If \(j-i>H\), all intervening witnesses are nonzero. The intervals of
\(H\) interior positions and their adjacent overlaps transmit a single
plane by the same rank argument. The first and last intervals of
\(H\) interior positions join the respective endpoint row to that plane
by the three-row determinant. The plane intersects \(R=0\) in a line
because it contains a row with \(R\ne0\). Thus the projections of the
two endpoints are again proportional. All zero-residue witnesses now
have proportional projections. Their first-to-last span is less
than \(d_1h<h\). Removing that span leaves a subwindow of length
\((s+6)h+1\) at one end of the original interval, with only nonzero
witnesses and with the \(G\)-spacing. In the absence of zero
witnesses one can choose any such subwindow. The preceding
contradiction completes the proof.
\end{proof}

\subsection{Exact digit targets and root templates}

\begin{theorem}\label{thm:selectable-counts}
Fix \(s\ge2\), \(C_0\ge0\) and \(0<\gamma<1/2\). In every window
of Lemma~\ref{lem:phase-residue} there is a run of \(\gg h\)
consecutive phases on which the following count holds. Put
\(A=B^{s-\gamma}\). For bounded fixed half-open intervals
\(I\subset\R\), \(J\subset[0,1)\) of positive length and every
\(\theta\in\R\), count the integers \(a\) such that \(a/A\in I\)
and the first \(sh\) digits of \(\{\theta+a b^m\alpha\}\) are
\(V^s\), with \(w=[V]_b\) and \(w/B\in J\). The count is
\begin{equation}\label{eq:exact-power-count}
        (1+o(1))\,|I||J|\,AB/M .
\end{equation}
It is uniform in \(\theta\) and in the selected phases. Restriction
to \(a\equiv a_0\bmod d\), with \(d\) fixed, divides the main term
by \(d\).

There are also constants \(\chi,\tau>0\) such that one may select
a run of \(\gg h\) consecutive phases in the same window for the
following count. Put \(l=\lfloor\chi h\rfloor\), \(E=b^l\), and
\(E_1=E\) or \(bE\). Prescribe the first \(\log_b E_1\) digits and
the last \(l\) digits of \(V\). Suppose the center of the prescribed
prefix cylinder lies in \([1/4,3/4]\), and take any
\(1\le d\le B^\tau\), \(a_0\bmod d\). The count in that congruence
class with the prescribed templates and \(a/A\in I\) is
\begin{equation}\label{eq:exact-template-count}
        (1+o(1))\,\frac{|I|A}{d L E_1 E}.
\end{equation}
The error is uniform in the stated template and congruence ranges
and in \(\theta\) and the selected phases.
\end{theorem}

\begin{proof}
For a real \(y\) let \(\langle y\rangle\in[-1/2,1/2)\) denote its
centered fractional part. Since \(M=1+LD\), the exact criterion
at \(x\in[0,1)\), for large \(h\), is
\begin{equation}\label{eq:target-centred}
 \lfloor Mx\rfloor=Lw\text{ for some }0\le w<B
 \quad\Longleftrightarrow\quad -x\le L\langle Dx\rangle<1-x.
\end{equation}
A hit gives \(Dx-w=(Mx-Lw-x)/L\), of magnitude less than \(1/L\).
Conversely take \(w\) to be the nearest integer to \(Dx\).
At a hit \(w=\lfloor Bx\rfloor\).

For the first assertion take small positive \(\lambda,\kappa,\rho\)
with \(2(\kappa+\gamma+\lambda+\rho)<1\) and \(\rho<1/16\).
Use Lemma~\ref{lem:phase-residue} with \(\gamma_0=\gamma+\lambda\),
and write \(\beta=b^m\alpha\), \(e(t)=\exp(2\pi i t)\).
Take smooth compactly supported \(g,\psi\) and a fixed
trigonometric polynomial \(P\). For samples \(x_a=\{\theta+a\beta\}\)
sum
\[
 g(a/A)P(x_a)\sum_{j\in\Z}\psi(L(Dx_a+j)).
\]
The frequencies are \(u+vD\), with \(u\) from the finite support
of \(P\) and \(v\in\Z\). The periodic sum has coefficients
\(L^{-1}\widehat\psi(v/L)\). With \(a=a_0+dn\), the nonzero
frequencies \(|v|\le L B^{\kappa/2}\) give a step
\(d(u+vD)\) in the range of \eqref{eq:phase-bad}, and
\[
 (A/d)\|d(u+vD)\beta\|\gg_d B^\lambda.
\]
The weighted sum over the progression has bound
\[
 \left|\sum_n g((a_0+dn)/A)e(nt)\right|
 \ll_N (A/d)\bigl(1+(A/d)\|t\|\bigr)^{-N}
\]
by finite differences. The coefficient absolute sum is bounded.
The errors inside and outside the cutoff are, respectively,
\(O((A/d)B^{-N\lambda})\) and
\(O((A/d)B^{-(N-1)\kappa/2})\). Take a fixed \(N\) large enough
that both are \(o(A/(dL))\). The zero frequency has asymptotic
\((A/(dL))\int g\int_0^1P\int\psi\). None of these bounds depends
on \(\theta\). Trigonometric approximation and smooth
majorants and minorants for rectangles give Lebesgue area in
\((x,t)\), \(t=L\langle Dx\rangle\), after division by \(A/(dL)\).
A finite union of rectangles can enclose and approximate the domain
\(x\in J,\ -x\le t<1-x\) up to arbitrarily small area, since the
boundary has area zero. Enlarging or shrinking \(J\) by a fixed small
amount and then sending that amount to zero handles \(w/B=x+O(1/B)\).
The same approximation on \(I\) yields \eqref{eq:exact-power-count}
since \(1/L\sim B/M\).

For the second assertion choose \(\chi,\tau,\mu,\lambda>0\) with
\[
 \gamma+4\chi+3\tau+4\mu+\lambda<1/2,\qquad
 2\chi+\tau+2\mu<1/16 .
\]
Set \(\kappa=\chi+\tau+2\mu,\ \rho=2\chi+\tau+2\mu\),
\(\gamma_0=\gamma+\chi+\tau+\lambda\). These satisfy the hypotheses
of Lemma~\ref{lem:phase-residue}. Let \(x_0\) be the prefix cylinder
center and \(w_0\bmod E\) the specified suffix label. Use the smooth
product
\[
  P_h(x)T_h(x)=\varphi(E_1(x-x_0))
      \sum_{t\in\Z}\psi\bigl(L(Dx-w_0-Et)\bigr)
\]
on \(0\le x<1\), extended periodically in \(x\). The first factor
is supported away from both endpoints for large \(h\). On the
support of the second factor, the integer from its nonzero term
must be the nearest integer \(w\) to \(Dx\), with \(w\equiv w_0\bmod E\).

Expand the second factor in \(Dx\), with period \(E\), and write
its index \(\nu=Ev+j\) for \(0\le j<E\). The factor
\(P_h(x)e(-jx/E)\) has Fourier coefficients \(p_j(u)\) with
\[
 |p_j(u)|\ll_N E_1^{-1}(1+|u|/E_1)^{-N}
\]
uniformly, by integrating derivatives on its support of length
\(O(1/E_1)\). The full coefficient is
\((LE)^{-1}p_j(u)\widehat\psi((v+j/E)/L)e(-jw_0/E)\)
at the integer frequency
\[
 k=u+vD+(B/E)j,\qquad E k=(Ev+j)D+(Eu+j).
\]
On \(|u|\le E_1B^\mu,\ |v|\le LB^\mu\) with \(d\le B^\tau\),
the integer \(dEk\) has magnitude at most \(B^{s+\kappa}\) and
residue representative \(d(Eu+j)\) of magnitude at most \(B^\rho<D\).
Thus \(k=0\) there forces \(u=v=j=0\). If \(k\ne0\), the inequality
\(\|dk\beta\|\le B^{-s+\gamma+\tau+\lambda}\) would make \(dEk\)
a forbidden frequency. Hence the step satisfies
\((A/d)\|dk\beta\|>B^\lambda\). The full coefficient absolute sum
is \(O(1)\), since \((LE)^{-1}\sum_{j,v}
|\widehat\psi((v+j/E)/L)|=O(1)\) and \(\sum_u|p_j(u)|=O(1)\).
The tail outside the cutoff has absolute sum
\(O(B^{-\mu(N-1)})\). As above, arbitrary fixed derivative order
makes each error \(o(A/(dLE_1E))\). The zero term equals
\((1+o(1))\,A(\int g)(\int\varphi)(\int\psi)/(dLE_1E)\)
for nonzero integrals.

Finally the prefix condition is \(E_1(x-x_0)\in[-1/2,1/2)\)
and on its slight enlargement \(x=x_0+O(1/E_1)\).
One can sandwich this condition, the interval \(I\), and the
\(t\)-interval \([-x,1-x)\) in \eqref{eq:target-centred} between
smooth functions with arbitrarily small integral discrepancy.
For the \(t\)-interval use translates of smooth approximants to
\([0,1)\), with supports modestly enlarged and shrunken by a
fixed margin to accommodate \(x-x_0\). All derivative bounds used
above are uniform for \(x_0\in[1/4,3/4]\). This gives
\eqref{eq:exact-template-count} with the specified half-open digit
conditions.
\end{proof}

\begin{corollary}\label{cor:phase-palindrome}
For \(s=2\), \(\theta=0\) and the phases of
Theorem~\ref{thm:selectable-counts}, one can choose
\(a_h\equiv1\bmod b\), \(0<a_h=o(B^{2-\gamma})\), such that the
actual digits of \(a_h\alpha\) after \(m\) places contain a square
of least period \(h\). At its seam there is a prescribed
palindrome \(\overleftarrow{U}U\) or \(\overleftarrow{U}cU\) of
length \(2l\) or \(2l+1\), respectively, provided the initial
template \(U\) or \(cU\) has its cylinder center in \([1/4,3/4]\).
Here \(c\) is a digit, \(l=\lfloor\chi h\rfloor\), and
\(\overleftarrow{U}\) is the reversal of \(U\).
\end{corollary}

\begin{proof}
Prescribe the indicated prefix and the suffix \(\overleftarrow{U}\).
For any fixed \(0<c_1<1\), the number of hits in \(0<a<c_1A\)
with \(a\equiv1\bmod b\) is \(\gg_{c_1} B^{1-\gamma-2\chi}\).
Distinct multipliers in this range have circle distance greater
than \(1/M\) by \eqref{eq:phase-ridout} with sufficiently small
\(\epsilon\). Each root label specifies a digit cylinder of
length \(1/M\), so the hits use distinct roots. There are
\(O(B^{1/2})\) roots of length \(h\) with a period at most \(h/2\).
If \(V^2\) has a period \(t<h\), comparing positions in the two
copies makes \(V\) invariant under cyclic shift by \(t\). It then
has the period \(\gcd(t,h)\le h/2\). The parameter choice makes
\(1-\gamma-2\chi>1/2\), so there are hits of least period \(h\)
for every \(c_1\). The minimum successful positive multiplier
is \(o(A)\). The seam has the required digits because they are
part of the two completed copies.
\end{proof}

